\documentclass[a4paper,12pt]{article}
\usepackage[utf8]{inputenc}
\usepackage[T1]{fontenc}
\usepackage[ngerman]{babel}
\usepackage{amsmath,amssymb}
\usepackage{booktabs}
\usepackage{enumitem}
\usepackage[margin=2.5cm]{geometry}
\usepackage{tikz}
\setlength{\emergencystretch}{2em}

\begin{document}

\begin{center}
{\LARGE\bfseries \"Ubungsblatt 7 -- L\"osungen}\\[0.5em]
{\large VC Logik}
\end{center}

\bigskip

%% ============================================================
%% AUFGABE 1
%% ============================================================
\subsection*{Aufgabe 1 \textnormal{(2 Punkte)} -- Vorw\"artsverkettung f\"ur $KB_1$}

\textbf{Was ist ein Literal?}
Ein Literal ist eine Variable oder die Negation einer Variable. $X$
hei\ss t \emph{positives} Literal, $\neg X$ hei\ss t \emph{negatives}
Literal. In der Klausel $\{X, \neg Y, \neg Z\}$ steht z.\,B.\ ein
positives Literal ($X$) und zwei negative ($\neg Y, \neg Z$).

\medskip
\textbf{Was bedeutet $\bot$?}
Das Symbol $\bot$ (gelesen ``Bottom'', das umgedrehte T) steht f\"ur die
Konstante \emph{falsch}. Eine Implikation $\dots \to \bot$ bedeutet:
sobald die Pr\"amisse wahr wird, w\"urde Falsch folgen -- die Pr\"amisse
darf also nie wahr werden. Wird $\bot$ aus einer Formel ableitbar, ist
die Formel \textbf{unerf\"ullbar}.

\medskip
\textbf{Was ist eine Horn-Klausel?}
Eine Klausel mit \emph{h\"ochstens einem positiven} Literal. Die Anzahl
positiver Literale entscheidet \"uber den Typ:
\begin{itemize}[nosep,leftmargin=2em]
\item \textbf{Fakt} (1 positives, 0 negative): $\{A\}$ entspricht der
      Aussage ``$A$ gilt''.
\item \textbf{Implikation} (1 positives, $\geq 1$ negative):
      $\{\neg A_1, \dots, \neg A_n, B\}$ entspricht
      $A_1 \wedge \dots \wedge A_n \to B$. Das positive Literal ist die
      Konklusion, die negierten werden zur Pr\"amisse.
\item \textbf{Constraint} (0 positive, nur negative -- auch
      ``Goal-Klausel''): $\{\neg A_1, \dots, \neg A_n\}$ entspricht
      $A_1 \wedge \dots \wedge A_n \to \bot$. Eine Constraint ist also
      ein \emph{Verbot}: die Pr\"amissen d\"urfen nie gleichzeitig wahr
      werden, sonst entsteht $\bot$.
\end{itemize}
Eine \emph{Horn-Formel} ist eine Konjunktion von Horn-Klauseln. Sobald
\emph{eine} Klausel zwei oder mehr positive Literale hat, ist die ganze
Formel keine Horn-Formel mehr.

\medskip
\textbf{Was ist Vorw\"artsverkettung?}
Verfahren f\"ur Horn-Formeln. Start mit den Fakten. Sobald alle Pr\"amissen
einer Implikation in der Faktenmenge sind, wird die Konklusion erg\"anzt.
F\"ur die \emph{Erf\"ullbarkeitspr\"ufung}: feuert eine Constraint-Klausel
$A_1 \wedge \dots \wedge A_n \to \bot$, wird $\bot$ abgeleitet -- damit ist
die Formel \textbf{unerf\"ullbar}. Wird kein Widerspruch erreicht, ist sie
erf\"ullbar.

\medskip
\textbf{Aufgabenstellung.}
\[
KB_1: \quad (P \vee S) \wedge (\neg P \vee \neg S) \wedge (\neg S \vee A) \wedge (\neg A \vee P) \wedge P \wedge S.
\]
$KB_1$ ist keine Horn-Formel; daher muss zuerst eine Horn-Variante
gefunden werden.

\bigskip
\textbf{Schritt 1: Klauseln klassifizieren.}

\begin{center}
\renewcommand{\arraystretch}{1.25}
\begin{tabular}{llcl}
\toprule
Klausel               & Positive Literale & Horn?                       & Bedeutung \\
\midrule
$\{P, S\}$            & $P$, $S$ \;(2)    & \textbf{nein} -- mehr als 1 & Disjunktion $P \vee S$ \\
$\{\neg P, \neg S\}$  & -- \;(0)          & ja                          & Constraint: $P \wedge S \to \bot$ \\
$\{\neg S, A\}$       & $A$ \;(1)         & ja                          & Implikation: $S \to A$ \\
$\{\neg A, P\}$       & $P$ \;(1)         & ja                          & Implikation: $A \to P$ \\
$\{P\}$               & $P$ \;(1)         & ja                          & Fakt: $P$ \\
$\{S\}$               & $S$ \;(1)         & ja                          & Fakt: $S$ \\
\bottomrule
\end{tabular}
\end{center}

\textbf{Warum $\{P, S\}$ st\"ort.}
Diese Klausel hat \emph{zwei} positive Literale ($P$ und $S$), eine
Horn-Klausel darf aber h\"ochstens \emph{eins} haben. Damit ist $KB_1$
keine Horn-Formel und Vorw\"artsverkettung kann nicht direkt angewendet
werden.

\bigskip
\textbf{Schritt 2: Horn-Variante finden.}

$\{P, S\}$ kann einfach \emph{weggelassen} werden, ohne dass sich an der
Erf\"ullbarkeit etwas \"andert. Die Begr\"undung in drei kleinen Schritten:
\begin{enumerate}[nosep,leftmargin=2em]
\item $\{P\}$ ist als eigene Klausel in $KB_1$ enthalten. In jedem Modell
      von $KB_1$ muss $P$ also wahr sein.
\item Wenn $P$ wahr ist, ist auch $P \vee S$ automatisch wahr (ein ODER
      wird wahr, sobald \emph{irgendein} Teil wahr ist -- $S$ spielt
      dann keine Rolle mehr).
\item Damit ist die Klausel $\{P, S\}$ \emph{redundant}: jede Belegung,
      die die anderen Klauseln erf\"ullt, macht ohnehin schon $P$ wahr
      und erf\"ullt damit auch $\{P, S\}$ ``gratis''.
\end{enumerate}
Eine redundante Klausel bringt keine zus\"atzliche Information und darf
gestrichen werden. \"Ubrig bleibt eine Formel, in der jede Klausel Horn
ist -- die gesuchte \textbf{Horn-Variante} $KB_1'$ (die untere Klammer
zeigt jeweils die Implikationsschreibweise der Klausel):
\[
KB_1' \;:\quad \underbrace{(\neg P \vee \neg S)}_{P \wedge S \to \bot}
\;\wedge\; \underbrace{(\neg S \vee A)}_{S \to A}
\;\wedge\; \underbrace{(\neg A \vee P)}_{A \to P}
\;\wedge\; P \;\wedge\; S.
\]
$KB_1$ ist genau dann erf\"ullbar, wenn $KB_1'$ erf\"ullbar ist -- die
Frage f\"ur $KB_1$ wird also durch die Antwort f\"ur $KB_1'$ vollst\"andig
beantwortet.

\bigskip
\textbf{Schritt 3: Vorw\"artsverkettung.}

Vorgehen kurz:
\begin{itemize}[nosep,leftmargin=2em]
\item \textbf{Schritt 0}: Fakten in die Faktenmenge schreiben.
\item \textbf{Weitere Schritte}: jede Regel, deren Pr\"amissen \emph{alle}
      schon Fakten sind, ``feuert'' -- ihre Konklusion wird neuer Fakt.
\item Feuert eine \textbf{Constraint} ($\dots \to \bot$), ist die Formel
      \textbf{unerf\"ullbar}.
\end{itemize}

\begin{center}
\small
\renewcommand{\arraystretch}{1.3}
\begin{tabular}{p{0.18\textwidth}p{0.30\textwidth}p{0.42\textwidth}}
\toprule
Schritt & Was passiert & Faktenmenge danach \\
\midrule
0: Start
& $\{P\}$ und $\{S\}$ direkt \"ubernehmen.
& $\{P, S\}$ \\
\midrule
1: $S \to A$
& Pr\"amisse $S$ erf\"ullt $\Rightarrow$ $A$ wird neuer Fakt.
& $\{P, S, A\}$ \\
\midrule
2: $A \to P$
& Pr\"amisse $A$ erf\"ullt $\Rightarrow$ $P$ w\"are neuer Fakt, ist aber schon dabei.
& $\{P, S, A\}$ \\
\midrule
3: $P \wedge S \to \bot$
& Pr\"amissen $P, S$ erf\"ullt $\Rightarrow$ $\bot$ wird abgeleitet.
& \textbf{Widerspruch}. \\
\bottomrule
\end{tabular}
\end{center}

In Schritt~3 feuert die Constraint und liefert $\bot$. Damit ist $KB_1'$
(und somit $KB_1$) \textbf{unerf\"ullbar}.

\medskip
\noindent\fbox{\parbox{\dimexpr\textwidth-2\fboxsep-2\fboxrule}{%
\textbf{Ergebnis Aufgabe 1.}\\[2pt]
$(P \vee S)$ ist redundant (folgt aus dem Fakt $P$) und kann weggelassen
werden -- das ergibt die Horn-Variante. Vorw\"artsverkettung leitet aus
den Fakten $P, S$ \"uber $S \to A$ und $A \to P$ nichts Neues ab, danach
feuert aber $P \wedge S \to \bot$. Damit ist $KB_1$ \textbf{unerf\"ullbar}.%
}}

\newpage

%% ============================================================
%% AUFGABE 2
%% ============================================================
\subsection*{Aufgabe 2 \textnormal{(2 Punkte)} -- Folgt $H$ aus $KB_2$?}

\textbf{Folgerung mittels Vorw\"artsverkettung.}
F\"ur Horn-Formeln gilt: $KB \models q$ genau dann, wenn $q$ durch
Vorw\"artsverkettung aus $KB$ ableitbar ist. Konkret: man h\"angt
neue Fakten an, bis $q$ erscheint (dann folgt $q$) oder keine neuen
Fakten mehr entstehen (dann folgt $q$ nicht).

\emph{In einfachen Worten:} Mit den gegebenen Fakten als Startpunkt so
viele neue Fakten wie m\"oglich ableiten. Taucht $H$ irgendwann in der
Faktenmenge auf, folgt $H$ aus $KB_2$. Wenn nichts Neues mehr abgeleitet
werden kann und $H$ nicht dabei ist, folgt $H$ nicht.

\medskip
\textbf{Aufgabenstellung.} $KB_2$:
\[
\begin{array}{ll}
R_1: A \wedge B \to K        & R_2: A \wedge B \wedge C \to M \\
R_3: C \wedge M \to D        & R_4: F \wedge M \to H \\
R_5: K \wedge D \wedge L \to E & R_6: L \wedge G \to F \\
R_7: A \wedge B \to L        & R_8: L \wedge M \to G \\
\multicolumn{2}{l}{\text{Fakt: } A \wedge B \wedge C}
\end{array}
\]
Folgt $H$?

\bigskip
\textbf{Vorw\"artsverkettung.}
Start: $\{A, B, C\}$. In jedem Schritt wird eine Implikation aktiviert,
deren Pr\"amissen vollst\"andig in der Faktenmenge sind.

\begin{center}
\small
\renewcommand{\arraystretch}{1.3}
\begin{tabular}{p{0.08\textwidth}p{0.30\textwidth}p{0.10\textwidth}p{0.40\textwidth}}
\toprule
Schritt & Aktivierte Regel & Neuer Fakt & Faktenmenge danach \\
\midrule
0 & -- (Start)                  & --   & $\{A, B, C\}$ \\
\midrule
1 & $R_1:\; A \wedge B \to K$   & $K$  & $\{A, B, C, K\}$ \\
\midrule
2 & $R_7:\; A \wedge B \to L$   & $L$  & $\{A, B, C, K, L\}$ \\
\midrule
3 & $R_2:\; A \wedge B \wedge C \to M$ & $M$ & $\{A, B, C, K, L, M\}$ \\
\midrule
4 & $R_3:\; C \wedge M \to D$   & $D$  & $\{A, B, C, K, L, M, D\}$ \\
\midrule
5 & $R_8:\; L \wedge M \to G$   & $G$  & $\{A, B, C, K, L, M, D, G\}$ \\
\midrule
6 & $R_6:\; L \wedge G \to F$   & $F$  & $\{A, B, C, K, L, M, D, G, F\}$ \\
\midrule
7 & $R_4:\; F \wedge M \to H$   & $\mathbf{H}$ & $\{A, B, C, K, L, M, D, G, F, \mathbf{H}\}$ \\
\bottomrule
\end{tabular}
\end{center}

$H$ erscheint in Schritt~7 in der Faktenmenge. (Zur Vollst\"andigkeit:
$R_5$ k\"onnte als 8.\ Schritt noch $E$ ableiten, das ist f\"ur die
Frage aber nicht n\"otig.)

\medskip
\noindent\fbox{\parbox{\dimexpr\textwidth-2\fboxsep-2\fboxrule}{%
\textbf{Ergebnis Aufgabe 2.}\\[2pt]
Vorw\"artsverkettung leitet aus dem Fakt $A \wedge B \wedge C$ \"uber die
Kette
\[
K,\; L,\; M,\; D,\; G,\; F \;\longrightarrow\; H
\]
das Atom $H$ ab. Damit gilt $KB_2 \models H$, also $H$ \textbf{folgt} aus
$KB_2$.%
}}

\newpage

%% ============================================================
%% AUFGABE 3
%% ============================================================
\subsection*{Aufgabe 3 \textnormal{(3 Punkte)} -- R\"uckw\"artsverkettung f\"ur $H$ aus $KB_3$}

\textbf{Vorgehen.} Erst $KB_3$ in Horn-Form bringen (jede Klausel mit
h\"ochstens einem positiven Literal als Implikation schreiben), dann
R\"uckw\"artsverkettung: Ziel $H$ rekursiv auf Teilziele zur\"uckf\"uhren,
bis nur Fakten \"ubrig sind.

\medskip
\textbf{Aufgabenstellung.} $KB_3$:
\[
\begin{array}{ll}
(H \wedge C) \vee (F \wedge A) \to M & A \wedge B \wedge C \\
B \wedge M \to J                     & \neg J \to \neg L \\
\neg B \vee \neg Y \vee L            & L \wedge G \to F \\
\neg(B \wedge C \wedge \neg G)       & \neg G \vee N \\
A \wedge J \to H                     & N \wedge F \to H \\
C \wedge N \to M
\end{array}
\]
Folgt $H$?

\bigskip
\textbf{Schritt 1: Horn-Form.}

\begin{center}
\renewcommand{\arraystretch}{1.25}
\begin{tabular}{ll}
\toprule
Original                                   & Horn-Form \\
\midrule
$(H \wedge C) \vee (F \wedge A) \to M$     & $H \wedge C \to M$ \;\textbf{und}\; $F \wedge A \to M$ \\
$\neg J \to \neg L$                        & $L \to J$ \;(Kontrapositiv) \\
$\neg B \vee \neg Y \vee L$                & $B \wedge Y \to L$ \\
$\neg(B \wedge C \wedge \neg G)$           & $\neg B \vee \neg C \vee G$ \;$\Rightarrow$\; $B \wedge C \to G$ \\
$\neg G \vee N$                            & $G \to N$ \\
\bottomrule
\end{tabular}
\end{center}

Damit ergibt sich die Horn-Regelmenge (Fakten: $A, B, C$):
\[
\begin{array}{rl@{\qquad}rl}
R_1: & H \wedge C \to M        & R_2: & F \wedge A \to M \\
R_3: & B \wedge M \to J        & R_4: & L \to J \\
R_5: & B \wedge Y \to L        & R_6: & L \wedge G \to F \\
R_7: & B \wedge C \to G        & R_8: & G \to N \\
R_9: & A \wedge J \to H        & R_{10}: & N \wedge F \to H \\
R_{11}: & C \wedge N \to M     & &
\end{array}
\]

\bigskip
\textbf{Schritt 2: R\"uckw\"artsverkettungs-Beweisbaum.}

\begin{center}
\begin{tikzpicture}[
  grow=down,
  level distance=1.2cm,
  level 1/.style={sibling distance=42mm},
  level 2/.style={sibling distance=42mm},
  level 3/.style={sibling distance=28mm},
  level 4/.style={sibling distance=22mm},
  level 5/.style={sibling distance=18mm},
  every node/.style={draw, rounded corners, inner sep=4pt, font=\small},
  edge from parent/.style={draw, -latex}
]
\node {$H$ \;via $R_9: A \wedge J \to H$}
  child {node {$A$ \checkmark\,Fakt}}
  child {node {$J$ \;via $R_3: B \wedge M \to J$}
    child {node {$B$ \checkmark\,Fakt}}
    child {node {$M$ \;via $R_{11}: C \wedge N \to M$}
      child {node {$C$ \checkmark\,Fakt}}
      child {node {$N$ \;via $R_8: G \to N$}
        child {node {$G$ \;via $R_7: B \wedge C \to G$}
          child {node {$B$ \checkmark\,Fakt}}
          child {node {$C$ \checkmark\,Fakt}}
        }
      }
    }
  };
\end{tikzpicture}
\end{center}

\medskip
\textbf{Sackgassen, die unterwegs aufgegeben werden.}
\begin{itemize}[leftmargin=2em,nosep]
\item Bei $M$ scheidet $R_1$ ($H \wedge C \to M$) aus, weil $H$ gerade
      das offene Hauptziel ist (Schleife).
\item Bei $M$ scheidet auch $R_2$ ($F \wedge A \to M$) aus: $F$ ginge
      nur \"uber $R_6: L \wedge G \to F$, und $L$ ginge nur \"uber
      $R_5: B \wedge Y \to L$ -- aber $Y$ ist weder Fakt noch
      ableitbar.
\end{itemize}
Daher der Weg \"uber $R_{11}: C \wedge N \to M$.

\medskip
\noindent\fbox{\parbox{\dimexpr\textwidth-2\fboxsep-2\fboxrule}{%
\textbf{Ergebnis Aufgabe 3.}\\[2pt]
R\"uckw\"artsverkettung findet den Beweisbaum
\[
H \;\leftarrow\; A,\, J \;\leftarrow\; B,\, M \;\leftarrow\; C,\, N \;\leftarrow\; G \;\leftarrow\; B,\, C,
\]
alle Bl\"atter sind Fakten. Damit \textbf{folgt} $H$ aus $KB_3$.%
}}

\end{document}
